\Article{14}
Nach dieser geometrischen Darstellung der Verhältnisse der horizontalen magnetischen Kraft entwickeln wir nun ein Verfahren zu ihrer Berechnung. Auf der Erdoberfläche ist $V$ nur noch eine Funktion zweier veränderlicher Größen. Dafür wählen wir die geografische Länge, die wir von einem beliebigen Ausgangsmeridian nach Osten zählen, und den Abstand vom Nordpol. Die Länge bezeichnen wir mit $\lambda$ und den Abstand, das Komplement der geografischen Breite, mit $u$. Betrachten wir die Erde als durch Drehung einer Ellipse mit großer Halbachse $=R$ und kleiner Halbachse $=(1-\varepsilon)R$ um die kleine Halbachse entstanden, so ist die Länge eines Meridianelements
\SourceLeaf{142}
\[
=\frac{(1-\varepsilon)^2R\cdot\mathrm{d}u}{(1-(2\varepsilon-\varepsilon\varepsilon)\cos u^2)^{\frac32}}
\]
und die Länge eines Elements des Breitenkreises
\[
=\frac{R\sin u\cdot\mathrm{d}\lambda}{\BareRoot{(1-(2\varepsilon-\varepsilon\varepsilon)\cos u^2)}}
\]

Zerlegt man die horizontale magnetische Kraft in zwei Teile, von denen der eine, $X$, in Richtung des Erdmeridians und der andere, $Y$, senkrecht dazu wirkt, und nimmt man $X$ als positiv, wenn diese Komponente nach Norden gerichtet ist, und $Y$ als positiv, wenn sie nach Westen gerichtet ist, so erhält man
\[
\begin{aligned}
X&=-\frac{(1-(2\varepsilon-\varepsilon\varepsilon)\cos u^2)^{\frac32}}{(1-\varepsilon)^2}\cdot\frac{\mathrm{d}V}{R\mathrm{d}u}\\
Y&=-\BareRoot{(1-(2\varepsilon-\varepsilon\varepsilon)\cos u^2)}\cdot\frac{\mathrm{d}V}{R\sin u\cdot\mathrm{d}\lambda}
\end{aligned}
\]
Die gesamte horizontale Kraft ist dann
\[
=\BareRoot{(XX+YY)}
\]
und der Tangens der Deklination
\[
=\frac YX
\]

Vernachlässigt man das Quadrat der Abplattung $\varepsilon$, so nehmen diese Ausdrücke die folgende Form an
\[
\begin{aligned}
X&=-(1+(2-3\cos u^2)\varepsilon)\cdot\frac{\mathrm{d}V}{R\mathrm{d}u}\\
Y&=-(1-\varepsilon\cos u^2)\cdot\frac{\mathrm{d}V}{R\sin u\cdot\mathrm{d}\lambda}
\end{aligned}
\]
oder, wenn man die Abplattung ganz außer Betracht lässt,
\[
X=-\frac{\mathrm{d}V}{R\mathrm{d}u}\qquad Y=-\frac{\mathrm{d}V}{R\sin u\cdot\mathrm{d}\lambda}
\]

Die bisher verfügbaren Beobachtungsdaten sind noch viel zu dürftig und größtenteils viel zu ungenau, als dass es schon ratsam wäre, die sphäroidische Gestalt der Erde zu berücksichtigen. Das wäre zwar an sich nicht schwierig, würde aber die Rechnungen ohne jeden Nutzen \SourceLeaf{143}erheblich erschweren. Wir bleiben daher hier bei den zuletzt angeführten Formeln und betrachten die Erde als Kugel mit dem Radius $=R$.

